\documentclass[12pt]{article}
\usepackage[utf8]{inputenc}
\usepackage[T1]{fontenc}
\usepackage{hyperref}
\title{Chung Language Resources for K-12 Curriculum Development}
\author{Compiled by Claude Sonnet 5.5 (claude-sonnet-5-5) for Andr\'as Kornai}
\date{}
\begin{document}
\maketitle

\section{Introduction}
Chung (ISO 639-3 scq) is the speakers' own name for the language that Khmer calls Sa'och. It is an Austroasiatic language of the Pearic group, spoken in Veal Renh village, Prey Nob District, Cambodia, with a related dialect in Kanchanaburi, Thailand. Estimates are about 20 speakers in Cambodia (2008) and about 10 in Thailand. Very little documentation exists, and no school materials were found. This dossier lists only resources that appeared in web searches.

\section{Primary Educational Resources}
None found. The language is not known to be taught in schools.

\section{Linguistic and Reference Materials}
\begin{itemize}
\item Sa'och language overview. \url{https://en.wikipedia.org/wiki/Sa%27och_language}
\item Pearic languages overview. \url{https://en.wikipedia.org/wiki/Pearic_languages}
\item Chong language overview, a related Pearic language. \url{https://en.wikipedia.org/wiki/Chong_language}
\item Glottolog reference record, which includes Pearic comparative material (the search listing did not name the work). \url{https://glottolog.org/resource/reference/id/153047}
\item Britannica entry on Pearic languages. \url{https://www.britannica.com/topic/Pearic-languages}
\end{itemize}
Robert Headley's 1978 English-Pearic vocabulary is cited in the search summary, but no direct link to it was found.

\section{Religious and Community Resources}
\begin{itemize}
\item Joshua Project page on Sa'och in Cambodia (people-group profile). \url{https://joshuaproject.net/people_groups/14754}
\end{itemize}
No scripture translation or community organisation was found.

\section{Digital and Audio Resources}
\begin{itemize}
\item Yale Southeast Asia Studies article on linguist Jean-Michel Filippi's effort to document the language. \url{https://seasia.yale.edu/node/831}
\item Endangered Languages Project entry. \url{https://endangeredlanguages.com/lang/1226}
\end{itemize}
No published audio corpus was confirmed.

\section{Curriculum Development Context}
With about ten to twenty speakers, mostly elderly and bilingual in Khmer, school curriculum work is not realistic at present. The practical priority is documentation: recordings, a word list and basic grammar, ideally in cooperation with the community and the researcher cited above. Any such material should use the name Chung, since speakers reject Sa'och as pejorative.
\end{document}